\chapter{Auto Loader, Incremental Ingestion, and Change Data Capture}
\index{Auto Loader}\index{cloudFiles}\index{incremental ingestion}\index{schema inference}%
\index{rescued data}\index{CDC}\index{SCD Type 2}\index{AUTO CDC}\index{checkpoint}%
\begin{objectives}
\begin{itemize}
    \item Explain Auto Loader's architecture: cloudFiles source, directory listing versus notification mode, and checkpointed exactly-once semantics.
    \item Configure schema inference, schema hints, evolution, and the rescued-data column.
    \item Build triggered incremental ingestion pipelines from cloud storage volumes.
    \item Apply change data capture with \texttt{MERGE INTO} and DLT \texttt{AUTO CDC}, including SCD Type 1 and Type 2.
    \item Handle out-of-order and duplicate CDC events with \texttt{SEQUENCE BY}.
    \item Design ingestion for high file rates plus nightly CDC without double-processing.
\end{itemize}
\end{objectives}

\begin{crosscloud}
Auto Loader is Databricks' \textbf{managed file-ingestion} answer; CDC application is the second
half of the pattern every platform implements.

\vspace{2mm}
{\footnotesize
\begin{tabular}{@{}p{3.0cm}p{3.4cm}p{3.2cm}p{3.2cm}@{}}
\toprule
\textbf{Capability} & \textbf{Databricks} & \textbf{AWS} & \textbf{Azure / GCP / Fabric} \\
\midrule
File ingestion & Auto Loader (cloudFiles) & S3 events + Glue bookmarks & Event Grid + ADF / GCS transfer + Dataflow / OneLake events \\
Exactly-once & Checkpoint + write-ahead offsets & Bookmarks / idempotent sinks & Watermarks / manifests \\
Schema drift & Inference + hints + evolution & Glue schema detection & ADF mapping drift / Dataflow \\
Malformed rows & \texttt{\_rescued\_data} column & Dead-letter patterns & Error outputs / DLQ \\
CDC source & Lakeflow Connect, partner tools & DMS CDC to S3 & ADF CDC / Datastream \\
CDC apply & \texttt{MERGE} / DLT \texttt{AUTO CDC} & Glue/EMR MERGE & T-SQL MERGE / BigQuery MERGE \\
\bottomrule
\end{tabular}}

\vspace{1mm}
\textbf{Snowflake} ingests with Snowpipe (notification-driven) and applies CDC with streams/tasks
or dynamic tables; \textbf{MongoDB} sources CDC naturally via change streams.
\end{crosscloud}

\section{Week 11 Roadmap and Mental Model}
Most production data arrives as files and change feeds, not clean API calls. Week 11 covers the two ingestion workhorses: \textbf{Auto Loader}, which turns a growing directory of files into an exactly-once incremental stream, and \textbf{CDC application}, which turns a log of inserts/updates/deletes into a current (SCD1) or historical (SCD2) table \cite{databricksAutoLoader,databricksCDC}.

The mental model: \textbf{ingestion is a state machine, and the checkpoint is the state}. Auto Loader's checkpoint records which files (or offsets) have been processed; the stream itself is conceptually stateless. Every exactly-once property, every recovery behavior, and every ``why did it reprocess everything'' mystery traces back to that checkpoint.

\begin{definitionbox}
\textbf{Auto Loader} (\texttt{cloudFiles}) is Databricks' incremental file-ingestion source: it discovers new files in cloud storage (by directory listing or cloud notifications), tracks processed files in a checkpoint, infers and evolves schemas, and rescues malformed records---delivering exactly-once ingestion with a single \texttt{.load()} call.
\end{definitionbox}

\begin{keyterms}
\textbf{Notification mode}: cloud-provider events (SQS/Event Grid/GCS) announce new files. \textbf{Directory listing}: periodic listing to detect files. \textbf{Schema hints}: user-supplied type overrides layered on inference. \textbf{\texttt{\_rescued\_data}}: the column catching unparseable records. \textbf{\texttt{SEQUENCE BY}}: the ordering key resolving out-of-order CDC events. \textbf{SCD}: slowly changing dimension---Type 1 overwrites, Type 2 versions history.
\end{keyterms}

\section{Monday: Auto Loader Architecture}
\subsection{Discovery: Listing versus Notification}
In \textbf{directory-listing} mode (default), Auto Loader lists the source path and diffs against the checkpoint---simple, and fine to millions of files with partitioned directory layouts. In \textbf{notification mode}, the cloud provider pushes file events to a queue Auto Loader subscribes to---required at very high file rates or deeply nested layouts where listing becomes the bottleneck \cite{databricksAutoLoader}.

\subsection{Exactly-Once and the Checkpoint}
Each micro-batch commits file offsets to the checkpoint before writing; on restart, processing resumes from the checkpoint. Files already committed are never reprocessed; files never committed are always processed. Consequences to internalize: the checkpoint is per-stream state---two streams cannot share one; deleting it reprocesses everything; moving it without care skips or duplicates.

\begin{pitfall}
\textbf{Deleting the checkpoint to ``fix'' a stuck stream.} The checkpoint is the memory of what has been ingested. Deleting it means a full reprocess---into a table that already has the data, unless you also reset the target. If reprocessing is truly the goal, plan it: snapshot or version the target, reset both sides, reconcile after.
\end{pitfall}

\section{Tuesday: Schemas and the Rescued Column}
\subsection{Inference, Hints, and Evolution}
Auto Loader samples files to infer a schema, persists it at \texttt{cloudFiles.schemaLocation}, and applies it consistently. \textbf{Schema hints} override inference for known-trouble columns (\texttt{cloudFiles.schemaHints}), e.g.\ IDs that must stay strings. Evolution: with \texttt{cloudFiles.schemaEvolutionMode} (or \texttt{mergeSchema} patterns downstream), new columns are absorbed deliberately rather than silently.

\subsection{The Rescued Data Column}
With \texttt{cloudFiles.rescuedDataColumn} set (default \texttt{\_rescued\_data} in Databricks runtimes), any record that fails parsing lands in that column as raw JSON instead of failing the batch. The pipeline survives; the malformed evidence is queryable; a steward process can reprocess or report.

\begin{lstlisting}[language=Python,caption={Production-shaped Auto Loader configuration.},label={lst:ch11-autoloader}]
stream = (spark.readStream.format("cloudFiles")
    .option("cloudFiles.format", "csv")
    .option("header", True)
    .option("cloudFiles.schemaHints", "order_id STRING, amount DECIMAL(12,2)")
    .option("cloudFiles.schemaLocation",
            "/Volumes/de_training/bronze/_schemas/orders")
    .option("cloudFiles.schemaEvolutionMode", "addNewColumns")
    .option("rescuedDataColumn", "_rescued_data")
    .load("/Volumes/de_training/bronze/landing/orders/"))

(stream.writeStream.format("delta")
    .option("checkpointLocation", "/Volumes/de_training/bronze/_ckpt/orders")
    .trigger(availableNow=True)          # batch economics on a stream source
    .toTable("de_training.bronze.orders_stream"))
\end{lstlisting}

\section{Wednesday: CDC with MERGE and AUTO CDC}
\subsection{The MERGE Pattern}
A CDC feed is a stream of \texttt{(key, operation, values, sequence)}. The hand-rolled application is the Chapter 2 MERGE with operation awareness and an ordering guard:

\begin{lstlisting}[language=SQL,caption={CDC application with delete handling and recency guard.},label={lst:ch11-cdc-merge}]
MERGE INTO de_training.silver.customers AS t
USING staging.cdc_batch AS s
ON t.customer_id = s.customer_id
WHEN MATCHED AND s.op = 'DELETE' THEN DELETE
WHEN MATCHED AND s.seq > t.seq THEN UPDATE SET *
WHEN NOT MATCHED AND s.op <> 'DELETE' THEN INSERT *;
\end{lstlisting}

\subsection{AUTO CDC for SCD}
DLT's \texttt{AUTO CDC} (formerly \texttt{apply\_changes}) packages this: declare the key, the sequence column, and the SCD type; the framework builds the correct stateful application \cite{databricksCDC}.

\begin{lstlisting}[language=Python,caption={SCD Type 2 customer history via AUTO CDC.},label={lst:ch11-autocdc}]
import dlt
from pyspark.sql import functions as F

dlt.create_streaming_table("silver_customers_scd2")

dlt.create_auto_cdc_flow(
    target="silver_customers_scd2",
    source="bronze_customer_cdc",
    keys=["customer_id"],
    sequence_by=F.col("updated_at"),     # out-of-order safety
    stored_as_scd_type=2,                # full history with __START_AT/__END_AT
    except_column_list=["op", "updated_at"])
\end{lstlisting}

\begin{definitionbox}
\textbf{SCD Type 1} overwrites: the table shows current truth only. \textbf{SCD Type 2} versions: each change closes the old row (\texttt{\_\_END\_AT}) and opens a new one, so ``what did we believe on March 3rd?'' is a query, not a restore. Choose Type 2 whenever downstream joins (features, finance allocations) depend on historical attribute values.
\end{definitionbox}

\section{Thursday Hands-on Project: Incremental Files plus SCD2 Customers}
Ingest daily CSV batches with triggered Auto Loader (including a mid-stream schema change and a corrupted row), then apply a CDC feed to a Silver customer table with full Type 2 history.
\index{hands-on lab}\index{Auto Loader!lab}\index{AUTO CDC!lab}

\begin{lstlisting}[language=Python,caption={Week 11 lab outline: ingestion drift drill + SCD2 verification.},label={lst:ch11-lab}]
# Part A: Auto Loader drift drill
#  1. Run the listing-mode stream from lst:ch11-autoloader (availableNow).
#  2. Drop batch 2 WITH a new column (loyalty_tier) -> absorbed by evolution.
#  3. Drop batch 3 with one malformed row -> lands in _rescued_data.
#  4. Query: SELECT _rescued_data FROM de_training.bronze.orders_stream
#           WHERE _rescued_data IS NOT NULL;

# Part B: CDC into SCD2 (DLT pipeline as in lst:ch11-autocdc)
#  5. Apply day 1 CDC (inserts + one update).
#  6. Apply day 2 CDC (another update + one delete), INCLUDING an
#     out-of-order event with an older updated_at -> must be ignored
#     by SEQUENCE BY.
#  7. Verify: the changed customer has exactly two rows, contiguous
#     __START_AT/__END_AT; the out-of-order event changed nothing.
\end{lstlisting}

\subsection*{Deliverables}
\begin{itemize}
    \item The Auto Loader stream definition with schema location, hints, evolution mode, and checkpoint---all under governed volumes.
    \item Evidence of the evolved column and the rescued row.
    \item The SCD2 table showing versioned customer history with contiguous validity windows.
\end{itemize}

\subsection*{Acceptance Criteria}
\begin{itemize}
    \item Batch 2's new column appears in Bronze without code edits; batch 3's bad row is rescued, not fatal.
    \item A same-checkpoint re-run ingests zero files (exactly-once proof).
    \item The out-of-order CDC event demonstrably does not corrupt the SCD2 history.
    \item The delete produces a closed interval (or Type 1 removal per the chosen flow), documented.
\end{itemize}

\subsection*{Stretch Goals}
\begin{itemize}
    \item Rebuild Part B with hand-written \texttt{MERGE} instead of \texttt{AUTO CDC} and compare the code you now maintain.
    \item Switch Auto Loader to notification mode and document the cloud wiring (queue, events, permissions).
    \item Write the steward query that reports rescued rows per day, and wire it as an expectation.
\end{itemize}

\section{Friday Use Case: 500 Files per Minute plus a Nightly CDC}
A payments platform receives partner settlement files at $\sim$500 small files/minute around the clock, and a nightly CDC extract from the core banking database. Both feed the same Silver domain. Design the ingestion layer.
\index{ingestion design}\index{payments}

\subsection{Requirements and Assumptions}
Settlement files are CSV, $\sim$2 KB each, arriving to a cloud container by partner/date hour. CDC is a nightly Debezium-exported feed. Downstream finance reconciliation runs at 06:00 and must reconcile to the penny against both sources.

\begin{figure}[H]
    \centering
    \resizebox{\textwidth}{!}{%
    \begin{tikzpicture}[
        src/.style={rectangle, rounded corners, draw=codegray, thick, fill=white, align=center, minimum width=3.0cm, minimum height=0.95cm, font=\small\bfseries},
        ing/.style={rectangle, rounded corners, draw=primary, thick, fill=primary!7, align=center, minimum width=3.2cm, minimum height=1.0cm, font=\small\bfseries},
        store/.style={cylinder, draw=primary, shape border rotate=90, aspect=0.25, fill=primary!10, align=center, minimum width=2.4cm, minimum height=1.15cm, font=\small\bfseries},
        flow/.style={-{Latex[length=2mm]}, draw=primary, thick}
    ]
        \node[src] (partners) at (0,1.2) {Partner files\\500/min};
        \node[src] (cdc) at (0,-1.2) {Core banking\\nightly CDC};
        \node[ing] (al) at (4.2,1.2) {Auto Loader\\notification mode\\triggered 5-min};
        \node[ing] (merge) at (4.2,-1.2) {AUTO CDC\\SCD2, SEQUENCE BY};
        \node[store] (bronze) at (8.4,1.2) {Bronze\\settlements};
        \node[store] (silver) at (8.4,-1.2) {Silver\\accounts SCD2};
        \node[ing, draw=orange!80!black, fill=orange!10] (recon) at (12.4,0) {06:00 reconciliation\\(Workflow gate)};
        \draw[flow] (partners) -- (al);
        \draw[flow] (cdc) -- (merge);
        \draw[flow] (al) -- (bronze);
        \draw[flow] (merge) -- (silver);
        \draw[flow] (bronze) -- (recon);
        \draw[flow] (silver) -- (recon);
    \end{tikzpicture}%
    }
    \caption{Dual-path ingestion: high-rate files via notification-mode Auto Loader, nightly CDC via SCD2, reconciled by a gated job.}
    \label{fig:ch11-ingestion}
\end{figure}

\subsection{Architecture Decisions}
\textbf{File path.} At 500 files/minute ($\sim$700k/day), directory listing degrades---notification mode (cloud events to a queue) keeps discovery O(new files). Triggered 5-minute runs bound latency and compact naturally (Chapter 5's small-file lesson applied at design time); the checkpoint and schema location live on governed volumes so the estate is auditable and redeployable.

\textbf{CDC path.} The nightly feed applies through \texttt{AUTO CDC} with \texttt{SEQUENCE BY updated\_at}; the accounts table is SCD2 because settlement disputes replay history. The 06:00 reconciliation is a Workflow task (Chapter 12) comparing Bronze settlement sums against Silver account movements---a failing gate pages the on-call, not the CFO.

\begin{pitfall}
\textbf{One stream, two sources.} Do not multiplex the CDC feed and the file stream into one checkpointed query ``for simplicity'': their rates, schemas, and recovery semantics differ. Two streams, two checkpoints, one reconciliation gate.
\end{pitfall}

\subsection{Risks and Mitigations}
\begin{table}[H]
\centering
\small
\begin{tabular}{@{}p{4.6cm}p{7.6cm}@{}}
\toprule
\textbf{Risk} & \textbf{Mitigation} \\
\midrule
Checkpoint loss double-ingests & Governed volumes, backup, and a documented reset-with-target procedure \\
Rescued data piles up unreviewed & Steward query weekly; spike alerts \\
Notification queue expires silently & Queue-depth and age-of-oldest-event alerts \\
CDC sequence column not unique & Uniqueness assertion on (\texttt{key}, \texttt{seq}) per batch \\
Schema hints rot as sources evolve & Hints versioned in the repo with the contract they encode \\
\bottomrule
\end{tabular}
\caption{Week 11 scenario risks and mitigations.}
\label{tab:ch11-risks}
\end{table}

\section{Design Decision Guide}
\index{decision guide!Week 11}%
Ingestion decisions set your recovery story; make them with the failure drill in mind.

\begin{table}[H]
\centering
\small
\begin{tabular}{@{}p{3.4cm}p{4.4cm}p{4.6cm}@{}}
\toprule
\textbf{Decision} & \textbf{Choose the first when} & \textbf{Choose the second when} \\
\midrule
Directory listing vs.\ notification mode & Moderate file counts, clean layouts & Millions of files or high arrival rates \\
Rescued data vs.\ fail-on-parse & Drifty external producers & Strict contracted feeds \\
SCD Type 1 vs.\ Type 2 & Current truth suffices & History feeds features, finance, or audit \\
\texttt{AUTO CDC} vs.\ hand-written MERGE & Standard upsert/SCD flows & Unusual sequencing or side effects \\
\texttt{availableNow} vs.\ processing-time & Batch economics on file arrival & Continuous minute-level freshness \\
Schema evolution vs.\ contract failure & Additive, tolerated change & Breaking change that must page someone \\
\bottomrule
\end{tabular}
\caption{Week 11 decision guide.}
\label{tab:ch11-decisions}
\end{table}

The two rules that override the table: one checkpoint per stream, and never delete state to fix a stall. Everything else is a judgment call you can revisit; those two are scars you cannot.

\section{Operational Guidance for Week 11}
\index{operational guidance!Week 11}%
\subsection{Performance and Reliability}
Checkpoints and schema locations are production assets: they live in governed volumes, appear in the runbook, and are covered by the same backup/DR thinking as tables. Steward the rescued-data column weekly---a spike means a producer changed shape. Trend files-ingested per run against files-arrived; divergence is a discovery lag or a stuck notification queue, not a data problem.

\subsection{Security and Governance Checklist}
\begin{itemize}
    \item Landing volumes readable only by the ingestion identity and stewards.
    \item Notification-mode queues and their credentials scoped to the source container.
    \item CDC feeds carry operation semantics under contract (\texttt{op}, sequence column documented).
    \item Rescued data containing PII is governed to the source's classification.
\end{itemize}

\subsection{Troubleshooting Playbook}
\begin{table}[H]
\centering
\small
\begin{tabular}{@{}p{3.6cm}p{4.2cm}p{4.8cm}@{}}
\toprule
\textbf{Symptom} & \textbf{Likely cause} & \textbf{First action} \\
\midrule
Stream finds no new files & Checkpoint ahead of arrivals / wrong path & Verify source path listing; check trigger ran \\
Everything reprocessed & Checkpoint lost or reset & Confirm target reset plan before re-running \\
Rescued-data spike & Producer schema drift & Diff a rescued row against the inferred schema \\
CDC applied out of order & Missing/wrong \texttt{SEQUENCE BY} & Verify the sequence column's monotonicity per key \\
\bottomrule
\end{tabular}
\caption{Week 11 troubleshooting quick reference.}
\label{tab:ch11-trouble}
\end{table}

\section{Interview Q\&A}
\subsection{Concepts and Trade-offs}
\begin{enumerate}
    \item \textbf{Q: How does Auto Loader achieve exactly-once?}\\
    A: The checkpoint durably records processed files/offsets per micro-batch commit; restarts resume from it. Exactly-once then depends on an idempotent sink---Delta commits---completing the contract.
    \item \textbf{Q: Notification mode versus directory listing?}\\
    A: Listing is simple and robust up to large file counts with clean layouts; notification mode scales to millions of files and high arrival rates by replacing listing with cloud events---at the cost of queue setup and permissions.
    \item \textbf{Q: What is the rescued-data column for?}\\
    A: Malformed records land there instead of failing the batch. The pipeline survives schema surprises; the evidence is queryable; reprocessing is a stewarded decision instead of a 3 a.m.\ incident.
    \item \textbf{Q: SCD1 or SCD2 for customer dimension?}\\
    A: SCD2 when downstream logic depends on historical attributes (features, finance allocations, dispute replay); SCD1 when only current truth matters and history is noise. When in doubt in regulated domains: SCD2.
    \item \textbf{Q: What does \texttt{SEQUENCE BY} protect against?}\\
    A: Out-of-order and replayed CDC events: without it, an older update can overwrite a newer one and silently regress your table state.
    \item \textbf{Q: How does Auto Loader differ from a scheduled \texttt{COPY} of new files?}\\
    A: Discovery, exactly-once tracking, schema management, and rescued-data handling are built in---a copy script rebuilds all four badly, usually discovering it during an incident.
    \item \textbf{Q: What is the first thing to check when a CDC table's history looks wrong?}\\
    A: The sequence column's uniqueness and monotonicity per key---nearly every SCD corruption is an ordering-contract violation upstream.
    \item \textbf{Q: What belongs in the ingestion contract with an external producer?}\\
    A: Schema and compatibility rules, key and sequence semantics, arrival SLA, and the malformed-data policy. Auto Loader tolerates drift; the contract decides which drift is tolerated versus paged.
\end{enumerate}

\section{Further Reading}
The Auto Loader documentation covers modes, schema management, and scaling \cite{databricksAutoLoader}; the CDC guide covers \texttt{AUTO CDC} flows and SCD types \cite{databricksCDC}. The underlying streaming semantics are in the Structured Streaming guide \cite{sparkStructuredStreaming}; Delta's MERGE and idempotency patterns trace to Chapter 2's foundations \cite{deltaLakeDocs}.

\section*{Key Takeaways}
\begin{itemize}
    \item Auto Loader = discovery + checkpoint + schema management; the checkpoint is the pipeline's memory---govern it.
    \item Notification mode is the scaling answer for high-rate file ingestion; listing is the simple default.
    \item Rescued data turns malformed records from incidents into stewarded backlog.
    \item CDC application is MERGE with operation awareness and a sequence guard; \texttt{AUTO CDC} packages it with SCD types.
    \item Separate sources get separate streams and checkpoints; reconciliation is an explicit, gated step.
\end{itemize}

\section{Exercises}
\begin{enumerate}
    \item \textbf{(Conceptual)} Why can't two streams share one checkpoint? What state would corrupt?
    \item \textbf{(Conceptual)} Explain the failure sequence if schema inference runs without a schema location set.
    \item \textbf{(Hands-on)} Complete the Thursday lab, including the out-of-order event test.
    \item \textbf{(Hands-on)} Build the steward report for rescued rows and wire it as a DLT expectation.
    \item \textbf{(Hands-on)} Replay the week's files into a recovery table using a fresh checkpoint; reconcile counts against the primary.
    \item \textbf{(Design)} Specify the notification-mode wiring (events, queue, permissions) for your cloud of choice.
    \item \textbf{(Design)} Design the 06:00 reconciliation gate for the Friday scenario: inputs, tolerances, failure action.
    \item \textbf{(Architecture)} At what arrival rate would you move the file path from triggered to continuous? Model the cost crossover.
\end{enumerate}
